% !TeX root = ../main.tex
\section{Fundamentos Teóricos y Trabajos Relacionados}
\label{sec:estado_arte}

La arquitectura propuesta converge tres pilares tecnológicos: registros distribuidos permisionados con consenso BFT, almacenamiento direccionable por contenido y aprendizaje federado. A continuación se sintetizan sus fundamentos teóricos y las limitaciones detectadas en la literatura.

\subsection{Blockchain Permisionada y Consenso IBFT 2.0}
Una cadena de bloques es una estructura de datos inmutable y distribuida en la que las transacciones se agrupan en bloques encadenados mediante funciones hash criptográficas resistentes a colisiones (como SHA-256)~\cite{nist2015sha, zheng2018blockchain}. Frente a las redes públicas (\textit{permissionless}) como Bitcoin o Ethereum Mainnet, penalizadas por altas latencias y costes económicos de transacción, los consorcios industriales adoptan redes permisionadas basadas en la Máquina Virtual de Ethereum (EVM)~\cite{buterin2014ethereum}, tales como Hyperledger Besu~\cite{besu_docs}.

En estos entornos, los mecanismos basados en consumo de recursos (PoW) o capital (PoS) se sustituyen por algoritmos de tolerancia a fallos bizantinos. El protocolo \textbf{IBFT 2.0} (\textit{Istanbul Byzantine Fault Tolerant})~\cite{moniz2020istanbul} estructura el consenso en torno a un conjunto predefinido de validadores que alternan el rol de proponente mediante un esquema rotatorio (\textit{round-robin}). Un bloque solo se anexa a la cadena tras una votación en tres fases (pre-prepare, prepare, commit) firmada por una supermayoría de al menos dos tercios:
\begin{equation}
n \ge 3f + 1
\end{equation}
donde $n$ denota el número total de nodos validadores y $f$ el número máximo de participantes defectuosos o bizantinos tolerables. 

IBFT 2.0 ofrece una propiedad determinante frente a otros consensos: la \textbf{finalidad determinista inmediata}. Al carecer de bifurcaciones probabilísticas (\textit{forks}), una vez que una ronda de entrenamiento o contribución se sella en un bloque, adquiere finalidad inmediata e inmutabilidad en el libro mayor distribuido. Asimismo, aunque el precio del gas (\textit{gas price}) se configure en cero para anular tarifas monetarias, el cómputo de \textit{gas used}~\cite{ethereum_gas} se mantiene como métrica objetiva del coste de cómputo y almacenamiento impuesto a la red.

\subsection{Almacenamiento Distribuido Off-Chain (IPFS)}
El almacenamiento directo de redes neuronales profundas en una Blockchain es inviable debido al tamaño de los pesos (decenas o cientos de megabytes), lo cual provocaría una saturación crítica del estado de la EVM (\textit{state bloat}) y un consumo prohibitivo de gas.

Para superar este cuello de botella, se implementa el paradigma de \textbf{almacenamiento fuera de cadena con compromiso dentro de cadena} (\textit{off-chain storage with on-chain commitment}) utilizando el Sistema de Archivos Interplanetario (IPFS)~\cite{benet2014ipfs, daniel2022ipfs}. IPFS opera sobre una topología P2P basada en direccionamiento por contenido (\textit{content addressing}) y estructuras de grafo acíclico dirigido (Merkle DAG)~\cite{daniel2022ipfs, trautwein2022design}: cada archivo se fragmenta en bloques inmutables y se identifica de forma unívoca mediante un \textit{Content Identifier} (CID), obtenido a partir del hash criptográfico de su contenido. La rigurosa evaluación arquitectural y empírica de IPFS a escala global realizada por Trautwein et al.~\cite{trautwein2022design} demuestra que su protocolo de intercambio de bloques (Bitswap) y la tabla hash distribuida (Kademlia DHT) posibilitan transferencias descentralizadas con baja sobrecarga de enrutamiento. Al desacoplar el almacenamiento de pesos neuronales en IPFS y registrar únicamente el identificador de contenido (CIDv0 de IPFS, que encapsula el multihash SHA-256 del artefacto serializado) en el contrato inteligente~\cite{zhao2021privacy}, se asegura que cualquier descarga posterior sea verificada contra el registro inmutable, facilitando deduplicación, verificación formal de integridad y acceso distribuido con coste computacional on-chain mínimo (debiendo asegurarse la disponibilidad y persistencia temporal o a largo plazo mediante anclaje de contenido o clústeres IPFS dedicados).

\subsection{Aprendizaje Federado y Algoritmo FedAvg}
El Aprendizaje Federado (FL)~\cite{mcmahan2017federated} resuelve la preservación de la privacidad de datos crudos al desacoplar el entrenamiento local de la agregación global. Entre los algoritmos de consenso federado, \textbf{FedAvg} (\textit{Federated Averaging}) representa el estándar fundacional. En cada ronda de comunicación $t$, el servidor distribuye el vector de pesos global actual $w^{(t)}$ a un subconjunto de $K$ clientes seleccionados. Cada cliente entrena el modelo sobre su partición local $\mathcal{D}_k$ minimizando una función de pérdida, obteniendo los parámetros locales $w_k^{(t)}$. El servidor calcula el nuevo modelo global consolidado como la media ponderada del volumen de muestras locales $n_k$:
\begin{equation}
w^{(t+1)} = \sum_{k=1}^{K} \frac{n_k}{n} \, w_k^{(t)}, \quad \text{donde } n = \sum_{k=1}^{K} n_k
\end{equation}

El marco de referencia \textbf{Flower}~\cite{beutel2020flower} proporciona una infraestructura desacoplada y eficiente para orquestar este flujo mediante llamadas a procedimientos remotos bidireccionales sobre gRPC.

\begin{table*}[t]
\centering
\caption{Comparativa cualitativa y arquitectural frente a los principales marcos de Blockchain-FL y robustez de la literatura. \\}
\label{tab:estado_arte}
\renewcommand{\arraystretch}{1.35}
\resizebox{\textwidth}{!}{%
\begin{tabular}{@{}l|ccccc@{}}
\toprule
\textbf{Propuesta / Framework} & \makecell{\textbf{Consenso}\\\textbf{Blockchain}} & \makecell{\textbf{Almacenamiento}\\\textbf{Modelo}} & \makecell{\textbf{Protocolo}\\\textbf{Verificación}} & \makecell{\textbf{Coste}\\\textbf{Transaccional}} & \makecell{\textbf{Tolerancia}\\\textbf{a Fallos}} \\
\midrule
\makecell[l]{Krum (Blanchard et al., 2017)~\cite{blanchard2017byzantine}} & \makecell{Sin Blockchain\\(Centralizado)} & \makecell{Servidor central\\(SPOF)} & \makecell{Selección geométrica\\($L_2$)} & \makecell{Inexistente\\(Sin ledger)} & \makecell{Bizantina estricta\\($f < \frac{N-2}{2}$)} \\ \midrule
\makecell[l]{BlockFL (Kim et al., 2019)~\cite{kim2019blockfl}} & \makecell{PoW\\(Lento / Ineficiente)} & \makecell{On-Chain\\(Cuello de botella)} & Unidireccional & \makecell{Criptomoneda\\volátil} & \makecell{Probabilística\\($< 50\%$)} \\ \midrule
\makecell[l]{DeepChain (Weng et al., 2021)~\cite{weng2021deepchain}} & \makecell{PoW / PBFT\\con incentivos} & \makecell{On-Chain\\(Gradients cifrados)} & \makecell{Cifrado Homomórfico\\+ MACs} & \makecell{Severo (Cifrado\\masivo on-chain)} & \makecell{Bizantina formal\\($f < N/3$)} \\ \midrule
AV-BFL (Pokhrel \& Choi, 2020)~\cite{pokhrel2020federated} & \makecell{PoW adaptativo /\\DPoS} & \makecell{On-Chain\\en RSUs} & \makecell{Verificación en\\estación base} & \makecell{Retransmisiones\\inalámbricas} & \makecell{Limitada por\\canal vehicular} \\ \midrule
\makecell[l]{H-BCFL (Chai et al., 2021)~\cite{chai2021hierarchical}} & \makecell{DPoS Jerárquico\\en RSUs} & \makecell{On-Chain en\\estaciones base} & \makecell{Reputación y\\auditoría de conocimiento} & \makecell{Sobrecarga por\\clustering en borde} & \makecell{Dependiente de\\fiabilidad de RSUs} \\ \midrule
B5G-FL (Lu et al., 2020)~\cite{lu2020blockchain} & \makecell{Permisionado\\(PBFT / BFT)} & \makecell{On-Chain en\\estaciones base} & \makecell{Consenso en\\estaciones base} & \makecell{Optimización por\\DRL} & \makecell{Bizantina formal\\($f < N/3$)} \\ \midrule
BFL Vehicular (Pujol et al., 2024)~\cite{pujol2024blockchain} & \makecell{IBFT 2.0\\(Permisionado)} & \makecell{Registro de\\alertas IoT} & \makecell{Monitoreo de\\tráfico} & \makecell{Libre de gas\\comercial} & \makecell{Bizantina formal\\($N \ge 3f + 1$)} \\ \midrule
\makecell[l]{BCFL-Ref (Goh et al., 2023)~\cite{goh2023blockchain}} & \makecell{PoW / PoA\\(Ethereum)} & \makecell{Híbrido Off-Chain\\(IPFS)} & \makecell{Unidireccional\\(Sin doble cruzada)} & \makecell{Gas variable\\(Red pública/privada)} & \makecell{Bizantina formal\\(Consenso Ethereum)} \\ \midrule
\makecell[l]{\textbf{Nuestra Arquitectura}} & \makecell{\textbf{QBFT / IBFT 2.0}\\\textbf{(Besu)}} & \makecell{\textbf{Almacenamiento}\\\textbf{por contenido (IPFS)}} & \makecell{\textbf{Doble Cruzada +}\\\textbf{Auditoría Ex-Post}} & \makecell{\textbf{Gas estable por}\\\textbf{transacción}} & \makecell{\textbf{Bizantina formal}\\\textbf{($f=1 \to 2$)}} \\
\bottomrule
\end{tabular}%
}
\end{table*}

\subsection{Vulnerabilidades en FL y Trabajos Relacionados}
A pesar de sus ventajas de privacidad, los sistemas tradicionales de Aprendizaje Federado sufren de vulnerabilidades intrínsecas clasificadas en recientes taxonomías y revisiones de seguridad~\cite{bouacida2021vulnerabilities, jere2021taxonomy, qu2022blockchain}:
\begin{enumerate}
    \item \textbf{Servidor central como punto único de fallo y confianza}: No existe una salvaguarda nativa para comprobar si el servidor agregador manipuló la ponderación de FedAvg o si inyectó sesgos maliciosos en $w^{(t+1)}$.
    \item \textbf{Ataques de envenenamiento y falta de no repudio}: Nodos clientes malintencionados pueden realizar envenenamiento de modelos (\textit{model poisoning}) o introducir puertas traseras (\textit{backdoors})~\cite{bagdasaryan2020backdoor}. Sin un historial no repudiable de identidades criptográficas y resúmenes de pesos, resulta inviable identificar y deslindar qué participante contaminó el modelo global consolidado.
\end{enumerate}

Para solventar estas carencias, la literatura especializada ha investigado la hibridación de Blockchain y FL (BCFL)~\cite{qu2022blockchain}. No obstante, el diseño de arquitecturas prácticas para sistemas vehiculares inteligentes (CAVs/IoV) plantea cuatro compromisos críticos:

\textbf{1. Almacenamiento y Cómputo On-Chain frente a Híbrido Off-Chain:} Las propuestas pioneras como BlockFL~\cite{kim2019blockfl} concibieron el intercambio completo de pesos neuronales a través del libro mayor distribuido. Asimismo, marcos como DeepChain~\cite{weng2021deepchain} articularon auditoría descentralizada e incentivos económicos combinados con cifrado homomórfico (HE) de gradientes, pero a costa de un severo cuello de botella computacional y de almacenamiento masivo on-chain. Otras alternativas como FLChain~\cite{bao2019flchain} introdujeron auditoría descentralizada mediante filtros de Bloom distribuidos, constatando que serializar matrices masivas de parámetros dentro de los bloques de la cadena satura el estado de la EVM (\textit{state bloat}) y degrada el rendimiento. Por ello, el paradigma híbrido (almacenamiento direccionable por contenido fuera de cadena mediante IPFS con anclaje criptográfico de hashes en cadena)~\cite{zhao2021privacy, goh2023blockchain} constituye una aproximación eficaz para mitigar la sobrecarga de estado y favorecer la escalabilidad en despliegues distribuidos.

\textbf{2. Mecanismos de Consenso y Latencia en Redes Inalámbricas y Vehiculares:} En la conducción autónoma y las redes de borde avanzadas (5G y más allá), los retardos temporales y la disponibilidad de recursos de cómputo y comunicación son factores determinantes~\cite{pokhrel2020federated, lu2020blockchain}. Los esquemas basados en Proof-of-Work (PoW) introducen demoras de varios minutos y bifurcaciones probabilísticas inasumibles para sistemas cooperativos en tiempo real~\cite{kim2019blockfl}. Para optimizar el compromiso entre consumo de recursos y convergencia del modelo en redes celulares de borde, Lu et al.~\cite{lu2020blockchain} formularon una arquitectura de Blockchain permisionada gestionada por estaciones base que emplea aprendizaje por refuerzo profundo (DRL) para la asignación dinámica de recursos, mientras que Pokhrel y Choi~\cite{pokhrel2020federated} modelaron analíticamente el impacto del canal inalámbrico vehicular en el retardo de consenso. En el dominio vehicular (IoV), Chai et al.~\cite{chai2021hierarchical} propusieron una arquitectura jerárquica delegando la agregación y el consenso DPoS en estaciones base en carretera (RSUs), Qi et al.~\cite{qi2021privacy} evaluaron una red permisionada combinada con privacidad diferencial para predicción de tráfico, y Abou El Houda et al.~\cite{abouelhouda2024blockchain} articularon esquemas de detección colaborativa de anomalías e intrusiones mediante BCFL en redes de borde vehiculares. No obstante, dichas aproximaciones condicionan la consistencia a la cobertura y sincronización de infraestructura vial fija. Frente a esquemas que asumen redes públicas con alta latencia y sobrecoste o alternativas que renuncian a la tolerancia a fallos bizantinos, los consorcios permisionados con consenso IBFT 2.0 / QBFT como Hyperledger Besu~\cite{pujol2024blockchain} proporcionan finalidad determinista inmediata y tolerancia bizantina rigurosa ($N \ge 3f + 1$) sin volatilidad económica de gas.

\textbf{3. Integridad Criptográfica frente a Robustez Estadística:} Conviene delimitar dos planos defensivos complementarios analizados en la literatura sobre seguridad en FL~\cite{bouacida2021vulnerabilities, jere2021taxonomy}. Por un lado, la capa de registro descentralizado e inmutable proporciona garantías formales sobre la integridad de los artefactos en tránsito, la autenticación de los emisores y el no repudio. Por otro lado, si un cliente manipula sus datos de entrenamiento de forma deliberada pero firma y registra legítimamente el resultado en la red (\textit{data/model poisoning}), se requiere incorporar algoritmos de agregación robusta frente a gradientes anómalos. La literatura ha propuesto esquemas de selección geométrica como Krum~\cite{blanchard2017byzantine}, filtros de eliminación por coordenadas como Trimmed-Mean y Median~\cite{yin2018byzantine}, o defensas híbridas frente a colusiones bizantinas complejas como Bulyan~\cite{mhamdi2018hidden}. De igual modo, marcos como FLTrust~\cite{cao2021fltrust} introdujeron el arranque de confianza (\textit{trust bootstrapping}) midiendo la similitud angular frente a un dataset de validación raíz en el servidor. Ambos planos resultan sinérgicos: la inmutabilidad y no repudio on-chain aíslan de forma inequívoca qué participante emitió cada vector de pesos, posibilitando la auditoría forense y la penalización (\textit{slashing}) coordinada desde el contrato inteligente ante la detección estadística de comportamientos disidentes.

\textbf{4. Fiscalización del Agregador (Caja Negra frente a Auditoría Ex-Post):} En la literatura sobre agregación verificable sin blockchain, propuestas pioneras como VerifyNet~\cite{xu2020verifynet} y VeriFL~\cite{guo2021verifl} han abordado la verificación criptográfica de la agregación mediante compromisos homomórficos, pruebas de conocimiento cero y esquemas de compartición de secretos. Si bien estos mecanismos garantizan formalmente que el servidor agrega correctamente sin aprender los gradientes individuales, imponen una sobrecarga computacional y de comunicación síncrona significativa sobre los nodos clientes en cada ronda federada. Por su parte, la mayoría de las arquitecturas BCFL asumen al servidor central como un componente fiduciario o intentan trasladar la agregación a la propia EVM con un coste computacional inasumible. Hasta donde alcanza nuestro conocimiento, los marcos existentes no proveen un mecanismo desacoplado y determinista que, sin penalizar el ciclo operativo en tiempo real ($<0{,}08\%$ de sobrecoste), faculte a cualquier participante o auditor externo a reconstruir \textit{ex-post} la agregación canónica de FedAvg a partir de los compromisos inmutables sellados en Blockchain e IPFS, verificando matemáticamente la fidelidad del modelo global sin vulnerar la privacidad de los datos locales.

Para contextualizar el aporte de esta investigación frente a las principales propuestas de la literatura especializada, la Tabla~\ref{tab:estado_arte} sintetiza una comparativa cualitativa y arquitectural. A diferencia de esquemas puramente estadísticos centralizados (como Krum), de arquitecturas basadas en consenso PoW o almacenamiento masivo on-chain (BlockFL, DeepChain), de propuestas dependientes de infraestructura jerárquica en RSUs o de modelos de verificación unidireccional sin doble cruzada (Goh et al.~\cite{goh2023blockchain}), nuestra arquitectura articula una doble verificación cruzada bidireccional junto con un protocolo de auditoría ex-post independiente sobre una red permisionada de consorcio con finalidad inmediata y coste de gas predecible y estable por transacción.


